\section{Validação e Consolidação do Modelo Integrado}

\subsection{Fundamentação Conceitual}

Como etapa conclusiva e integrada da metodologia, esta seção sintetiza e consolida os elementos delineados nas seções anteriores, transformando uma arquitetura conceitual sofisticada e um fluxo de processos teoricamente bem fundamentado em mecanismos práticos, operacionalizáveis e sustentáveis de gestão de conhecimentos tradicionais e de tomada de decisão. Para atender ao objetivo específico de criação, teste rigoroso e validação do modelo integrado (OE6), será aplicada uma metodologia de pesquisa-ação participativa conforme descrita em \cite{Tidd2005}, operacionalizada mediante uma abordagem sistêmica baseada em ciclos iterativos de planejamento colaborativo, ação orientada por objetivos, observação sistemática e reflexão coletiva sobre o que funcionou, o que não funcionou e por quê.

Esta abordagem não é meramente uma questão de preferência metodológica. Ela emerge da convicção de que modelos para gestão de conhecimentos tradicionais que não passam pelo crivo da validação participativa com as comunidades cujos saberes estão sendo estudados correm risco elevado de fracasso operacional ou, pior, de capturar valor sem entregar benefícios reais às comunidades. Portanto, a validação que se segue é simultaneamente um teste de adequação técnica e um processo de negociação e ajuste social do modelo.

\subsection{Arquitetura Operacional do Modelo Integrado} 

A consolidação do modelo integrado repousa em dois eixos fundamentais que funcionam de forma interdependente e mutuamente reforçadora. O primeiro eixo compreende os quatro módulos tecnológicos e analíticos já descritos em detalhe nas seções anteriores: o Módulo de Gestão de Propriedade Intelectual e Marcos Legais, o Módulo de Valoração de Ativos Intangíveis, o Módulo de Mapeamento e Classificação via Machine Learning, e o Módulo de Otimização e Tomada de Decisão. Estes módulos operam em profunda sinergia para transformar conhecimentos tradicionais mapeados e documentados em ativos intangíveis formalmente protegidos por mecanismos jurídicos adequados, economicamente valorados conforme múltiplas dimensões de valor, e continuamente priorizados e otimizados em cenários que levam em conta potenciais conflitos entre objetivos.

O segundo eixo é estruturalmente tão importante quanto o primeiro: é o eixo de governança colaborativa e democrática. Este eixo é operacionalizado através de três mecanismos complementares já descritos: a plataforma digital integrada que funciona como infraestrutura de conhecimento compartilhado e de tomada de decisão suportada por dados; o Comitê de Integração Metodológica e Governança Colaborativa que reúne representantes heterogêneos e exerce função de deliberação e validação; e os processos ágeis e iterativos de cocriação que garantem participação contínua de comunidades e refinamento permanente de ferramentas e processos conforme feedback real.

O resultado da articulação destes dois eixos é um modelo que não apenas realiza análises técnicas rigorosas e matematicamente sofisticadas (análise de marcos legais, valoração econômica de intangíveis, machine learning, otimização multiobjetivo), mas que o faz de forma estruturalmente participativa, respeitando plenamente autonomia e agência comunitária, validando conhecimentos tradicionais como contribuições científicas legítimas, e operando sob marcos legais e éticos que reconhecem direitos inalienáveis de povos e comunidades tradicionais. A validação deste modelo integrado, tal como será descrita a seguir, ocorrerá mediante uma abordagem multinível que combina rigor científico com legitimidade social, permitindo identificar simultaneamente se o modelo funciona tecnicamente e se as comunidades o reconhecem como legítimo e útil.

\subsection{Validação Participativa Multinível}

A validação do modelo integrado será conduzida mediante uma abordagem multinível sofisticada que combina simultaneamente rigor científico com legitimidade social, aplicando triangulação metodológica e pluralismo epistemológico para garantir que o modelo não apenas é coerente do ponto de vista técnico, mas é também percebido como apropriado, justo e útil pelos diferentes grupos de atores que o utilizarão.

\subsubsection{Nível 1 — Validação Técnico-Científica}

A primeira camada de validação aborda a adequação metodológica e o rigor técnico de cada componente. Esta etapa incluirá avaliação por pares de alto nível — convocando-se especialistas reconhecidos nas áreas de propriedade intelectual e direitos de povos tradicionais, valoração de ativos intangíveis, aprendizagem de máquina aplicada e otimização multiobjetivo. Os avaliadores examinarão relatórios técnicos especializados de cada módulo, focando em aspectos como adequação metodológica (o método escolhido é apropriado para a questão de pesquisa?), consistência interna (as conclusões se seguem logicamente das evidências?), alinhamento com o estado da arte (a pesquisa incorpora os avanços mais recentes em sua área?), e potencial de replicabilidade (outros pesquisadores poderiam reproduzir os resultados e expandir a metodologia?).

Em paralelo, será conduzido benchmarking internacional sistemático, comparando-se os resultados e metodologias com frameworks e boas práticas internacionalmente reconhecidas. Particularmente relevantes serão: as bases de dados de Domínio Público e Conhecimento Tradicional da Convenção sobre Diversidade Biológica, a Traditional Knowledge Digital Library mantida pela WIPO em parceria com comunidades indígenas globais, e os sistemas operacionalizados de Indicações Geográficas pelo INPI no Brasil. Este benchmarking permitirá identificar não apenas lacunas remanescentes, mas também contribuições originais que o projeto traz relativamente ao que já existe na literatura e prática internacional.

Para os componentes computacionais e de machine learning, serão realizados testes de desempenho rigorosos contemplando múltiplas dimensões: acurácia preditiva será avaliada via validação cruzada de k-fold e teste de holdout com comparação a baselines robustos da literatura (o modelo proposto supera significativamente algoritmos padrão?); estabilidade será investigada através de sensibilidade a variações nos dados de entrada (o modelo é robusto a outliers, dados faltantes, variações?) e sensibilidade a variações nos parâmetros (pequenas mudanças de parâmetros levam a mudanças drásticas de output?); escalabilidade será testada com diferentes volumes de dados para confirmar que tempos de processamento permanecem viáveis; e interpretabilidade será avaliada usando métodos formais como SHAP (SHapley Additive exPlanations) e LIME (Local Interpretable Model-agnostic Explanations), garantindo que as decisões do modelo podem ser explicadas, não apenas que produz previsões corretas.

\subsubsection{Nível 2 — Validação Operacional mediante Aplicação Piloto}

A segunda camada de validação leva o modelo integralmente para fora do laboratório/escritório de pesquisadores e o testa em contextos reais com comunidades. A aplicação piloto contemplará a seleção intencional de três a cinco comunidades quilombolas do Território de Identidade Semiárido Nordeste II, sendo que esta seleção não será aleatória mas orientada por critérios específicos: buscam-se comunidades que apresentem perfis distintos em termos de estrutura de conhecimento (algumas com saberes mais agrícolas, outras com saberes de extrativismo, outras com saberes de medicina tradicional), nível de engajamento institucional anterior (desde comunidades com experiências prévias de parceria com universidades até comunidades inteiramente novas nisto), e acesso a tecnologia (desde comunidades com internet adequada até comunidades com acesso limitado). Esta diversidade de contextos permite testar a robustez e adaptabilidade do modelo.

A implementação será cuidadosamente assistida por equipe dedicada de pesquisa que viverá junto às comunidades durante fases intensivas do projeto. As atividades de campo incluirão: oficinas de apresentação clara e acessível do modelo, não com jargão técnico mas com linguagem e exemplos relevantes localmente; processos participativos de capacitação em cada um dos módulos (como mapear conhecimentos, como valorá-los, como protegê-los); mapeamento participativo propriamente dito de conhecimentos tradicionais, conduzido conjuntamente com mestres de saberes e associações comunitárias; conduções de valoração multidimensional com grupos comunitários para avaliar concordância/discordância com resultados técnicos; estruturação participativa de estratégias de proteção de PI conforme prioridades comunitárias; e exploração conjunta de cenários de decisão para entender quais trade-offs entre objetivos faz sentido para as comunidades.

O processo será documentado sistematicamente não apenas quanto a resultados (quantos conhecimentos foram mapeados, quantos foram protegidos), mas em detalhe quanto a desafios encontrados (qual aspecto do modelo era difícil de entender? Que adaptações técnicas foram necessárias? Quais barreiras culturais ou institucionais emergiram?), e quanto a adaptações realizadas. Esta documentação não é um apêndice, mas parte central da validação, pois permite aprender sobre a transferibilidade do modelo.

Resultados intermediários serão monitorados por período de 12-18 meses pós-implementação, para avaliar sustentabilidade. Indicadores incluem: número de conhecimentos efetivamente documentados, catalogados, valorados e, idealmente, protegidos; parcerias firmadas entre comunidades e instituições externas (universidades, empresas, órgãos ambientais) que resultaram em acordos concretos; variações em indicadores de agência comunitária (as comunidades sentiram-se empoderadas para tomar decisões sobre seus conhecimentos?); e preservação ou ampliação de práticas culturais associadas aos conhecimentos (o modelo incentivou ou desestimulou transmissão intergeracional?).

\subsubsection{Nível 3 — Validação Social e Cultural (Legitimidade Comunitária e Representativa)}

A terceira camada de validação desloca o foco de se o modelo funciona para se o modelo é legítimo, apropriado e sustentável do ponto de vista das comunidades e da sociedade em sentido amplo. Concluída a fase piloto, serão realizadas oficinas estruturadas de devolutiva e coavaliação nas comunidades participantes. Estas oficinas não apresentam resultados de forma unidirecional (pesquisadores presenteando dados aos comunitários), mas funcionam como espaços de negociação de significados. Apresentam-se os resultados em linguagem acessível (com gráficos, fotos, narrativas, não apenas tabelas numéricas), e facilita-se discussão estruturada sobre: Foi o modelo útil? Que limitações identificaram? Quais barreiras enfrentaram na sua utilização? O modelo respeitou adequadamente conhecimentos sensíveis e estruturas de decisão tradicionais? É sustentável para as comunidades continuar utilizando sem dependência permanente de pesquisadores externos?

O processo observará rigor ético inquestionável, incluindo aprovação prévia por Comitê de Ética em Pesquisa, obtenção de consentimento livre, prévio e informado (individual para alguns participantes, coletivo para decisões de comunidade), garantia de confidencialidade conforme solicitado (alguns conhecimentos a comunidade deseja permaneçam privados), direito inquestionável de retirada sem prejuízo, repartição equitativa de benefícios (publicações, reconhecimento, recursos financeiros), e respeito total a conhecimentos identificados como sensíveis ou sagrados — com mecanismos de controle de acesso que a comunidade controla completamente.

Os resultados e materiais serão igualmente apresentados a organizações representativas de povos e comunidades em nível estadual e nacional — a Coordenação Estadual das Comunidades Quilombolas de Sergipe, a Coordenação Nacional de Articulação das Comunidades Negras Rurais Quilombolas, o Conselho Nacional de Povos e Comunidades Tradicionais — com vistas a obter validação social ampliada, avaliar potencial de replicação em outros contextos, e recolher sugestões de aperfeiçoamento que surjam da experiência cumulativa destes atores.

\subsection{Indicadores de Sucesso, Avaliação de Impacto e Teoria da Mudança}

A efetividade do modelo integrado será avaliada mediante um framework abrangente de indicadores estruturado explicitamente em teoria da mudança (Theory of Change). Conforme este marco conceitual, uma intervenção produz impactos apenas se há uma cadeia causal clara: investimentos (\textit{inputs}) são utilizados para realizar atividades específicas, que produzem \textit{outputs} tangíveis (produtos, serviços), que por sua vez geram \textit{outcomes} (mudanças no comportamento, atitude, capacidade de atores), que eventualmente convergem em impactos (mudanças estruturais, sistêmicas de longo prazo). Cada uma destas etapas requer supostos sobre como a mudança ocorre (por exemplo, a suposição que se comunidades tiverem conhecimento sobre proteção de PI, utilizarão essa proteção; ou que se conhecimentos forem valorados, sua preservação será priorizada).

Os indicadores foram organizados em três dimensões complementares que capturam diferentes aspectos do modelo e seus efeitos.

\subsubsection{Dimensão 1 — Gestão do Conhecimento, Proteção Jurídica e Apropriação de Valor}

Nesta dimensão, monitoram-se indicadores relacionados aos produtos diretos do modelo e ao acesso de comunidades a mecanismos formais de proteção. Indicadores quantitativos incluem: número absoluto de conhecimentos tradicionais documentados e catalogados em base de dados estruturada (com meta de, por exemplo, 200+ conhecimentos por comunidade piloto); percentual de conhecimentos para os quais foi identificado valor estratégico elevado via algoritmos de machine learning; quantidade e tipos de pedidos de proteção de propriedade intelectual efetivamente depositados (quantas marcas coletivas, quantas indicações geográficas, quantos registros defensivos); número e valor de acordos de licenciamento firmados entre comunidades e parceiros externos; e médias de variação no tempo de tramitação de processos administrativos (time-to-protect) comparando baseline pré-modelo com pós-implementação.

Além de quantidades, monitoram-se resultados qualitativos e processuais: taxa de adoção efetiva de protocolos de Consulta Livre, Prévia e Informada — não apenas que protocolos existem, mas que comunidades efetivamente os utilizam em decisões; ocorrências documentadas de prevenção de apropriação indevida ou biopirataria (por exemplo, quando vigilância via bases de dados defensivas identifica pedidos de patente inadequados que puderam ser contestados); e percepção de comunidades quanto à equidade e transparência nos processos de proteção (entrevistas qualitativas sobre se sentiram ouvidas, se custos foram considerados acessíveis, se cronogramas permitiram deliberação adequada).

\subsubsection{Dimensão 2 — Valoração Multidimensional e Priori­zação}

Nesta dimensão, avaliam-se resultados dos módulos de valoração de ativos. Indicadores incluem: volume bruto de conhecimentos submetidos a processos de valoração psicométrica e econômica (meta de 100+ conhecimentos por comunidade); índices de valor cultural, científico e econômico gerados para cada conhecimento (permitindo comparação, por exemplo, entre conhecimentos com alto valor cultural mas baixo valor econômico versus conhecimentos com ambos os valores elevados); e matrizes de priorização geradas que incorporam múltiplos critérios.

Um indicador particularmente relevante é a relação entre valoração técnica (aquela derivada de algoritmos, dados quantitativos) e percepção comunitária de importância (derivada de consulta, entrevistas com mestres de saberes). Espera-se alguma correlação, mas desvios são informativos: conhecimentos que a comunidade valoriza culturalmente mas o modelo classifica com baixo potencial comercial devem permanecer protegidos mesmo que não rentem economicamente. Adicionalmente, rastreia-se a proporção de conhecimentos priorizados pelo modelo que posteriormente se converteram em oportunidades de impacto territorial concreto (parcerias, pesquisas, iniciativas de comercialização) — se o modelo não identificar conhecimentos que geram impacto real, sua utilidade é questionável.

\subsubsection{Dimensão 3 — Sustentabilidade Cultural, Socioambiental e Bem-Estar Comunitário}

Esta dimensão é talvez a mais importante eticamente, pois avalia se o modelo contribui para objetivos de longo prazo que vão além de indicadores econômicos. Serão monitorados produtos e ações de valorização cultural originários do projeto — quantos documentários, publicações, exposições, eventos comunitários aconteceram? Qual o alcance (quantas pessoas alcançadas, quantas comunidades envolvidas)? Qual o engajamento de jovens em atividades de transmissão de conhecimentos — o modelo incentivou gerações mais jovens a aprender de mestres? Quantas novas parcerias educacionais (escolas, universidades) foram estimuladas?

Indicadores de resultado de longo prazo incluem: variações na percepção de identidade cultural entre participantes (avaliadas via surveys pré e pós, escalas psicométricas de identidade e enraizamento comunitário); taxa de transmissão intergeracional de práticas (porcentagem de jovens que aprenderam e praticam conhecimentos tradicionais); manutenção ou ampliação de áreas manejadas conforme práticas tradicionais (especialmente relevante para conhecimentos agroecológicos); e métricas de agrobiodiversidade e saúde ecossistêmica (variação em riqueza de espécies, resiliência a secas, etc., em áreas manejadas tradicionalmente versus áreas sob manejo convencional).

\subsubsection{Métodos de Coleta de Dados}

A coleta de indicadores utilizará pluralidade metodológica apropriada a cada tipo. Pesquisas estruturadas (questionários) permitirão coletar medidas de baseline antes da implementação e medidas de acompanhamento após, para aqueles indicadores que são passíveis de quantificação. Análise documental de contratos, registros administrativos, pedidos de proteção de PI, fornecerá dados sobre outputs formais. Entrevistas semiestruturadas e grupos focais com participantes-chave (mestres de saberes, lideranças, jovens) permitirão explorar de forma nuançada como o modelo foi percebido, que benefícios foram tangidos, quais expectativas não foram atendidas. Observação participante durante a implementação gerará notas de campo ricas sobre dinâmicas de engajamento, conflitos, adaptações. Registros fotográficos e audiovisuais documentarão atividades, celebrações, produtos comunitários. E consulta a bases secundárias (INPI para dados sobre proteções de PI, IPHAN para patrimônio, IBGE para dados sociodemográficos, sistemas de monitoramento de unidades de conservação) permitirá triangulação com fontes independentes.

\subsubsection{Análise de Dados e Inferência Causal}

A análise combinará rigorosamente métodos quantitativos com métodos qualitativos. Análise quantitativa incluirá estatística descritiva (médias, distribuições), testes de significância para detectar se variações pré-pós são genuínas (testes T, Mann-Whitney conforme distribuição), correlações entre indicadores, análises de regressão para explorar relações entre variáveis. Quando possível — particularmente se múltiplas comunidades participarem com temporalidade defasada — tentaremos aplicar métodos mais robustos de avaliação de impacto como desenho de diferenças em diferenças (difference-in-differences) comparando comunidades participantes com comunidades de controle, ou estudos de caso comparativos que pareiem comunidades similares e controlam para fatores contextuais.

Análise qualitativa incluirá codificação temática de entrevistas e grupos focais para identificar padrões em percepções e experiências, análise de conteúdo de narrativas comunitárias para extrair significados, e triangulação entre múltiplas fontes e métodos para validação cruzada de inferências. Um resultado é especialmente relevante: identificar não apenas se mudanças ocorreram, mas se as comunidades atribuem essas mudanças ao modelo (isso seria evidência de impacto causal) ou se reconhecem que outras forças (contexto político, condições de mercado, eventos aleatórios) influenciaram igualmente ou mais.

\subsection{Disseminação, Transferibilidade, Sustentabilidade Institucional e Perpetuação de Impactos}

\subsubsection{Produtos, Publicações e Saídas de Conhecimento}

Os resultados abrangentes da validação e do projeto como um todo serão sistematizados e disseminados através de múltiplos canais e formatos apropriados a públicos distintos.

Um \textbf{Relatório Técnico-Científico Abrangente} documentará integralmente a metodologia empregada (justificativas epistemológicas, escolhas técnicas), os resultados obtidos em cada módulo (com estatísticas, tabelas, gráficos), as análises de integração (como os módulos trabalharam conjuntamente, quais sinergias emergiram), lições aprendidas (o que funcionou bem, o que precisaria ser melhorado, por quê), e recomendações para futuros utilizadores do modelo. Este relatório será produzido em versão acadêmica formal para depositário em repositório institucional da universidade, mas também em versão acessível (linguagem clara, resumos em português simples) para disponibilização em plataformas abertas.

Um \textbf{Manual Operacional e Guia de Replicação} fornecerá aos potenciais utilizadores (pesquisadores, gestores ambientais, facilitadores comunitários) instruções passo-a-passo de como implementar o modelo em novos contextos. Este manual incluirá linguagem acessível complementada por fluxogramas visuais, listas de verificação de atividades, guias de tomada de decisão (se você se deparar com isso, então faça aquilo), e estudos de caso ilustrando diferentes cenários (comunidade com alta escolaridade vs. baixa, comunidade com tecnologia disponível vs. limitada, etc.). O manual também incluirá dicas práticas sobre gestão de stakeholders, resolução de conflitos, adaptação cultural de ferramentas.

\textbf{Artigos científicos e capítulos de livros} serão submetidos a periódicos de alto impacto nas múltiplas áreas envolvidas: propriedade intelectual e direitos de povos tradicionais (revistas em direito, gestão pública, estudos de povos tradicionais); economia da inovação e valoração de intangíveis (revistas de economia, administração, empreendedorismo); machine learning aplicado (revistas de inteligência artificial, processamento de dados); e otimização multiobjetivo e teoria da decisão (revistas de pesquisa operacional, planejamento estratégico). Esta estratégia multidisciplinar de publicação reflete a natureza verdadeiramente interdisciplinar do projeto e maximize o alcance de resultados a diferentes comunidades científicas.

\textbf{Notas Técnicas e Posicionamentos para Gestores Públicos} serão produzidas em linguagem executiva apropriada para formuladores de política nos diferentes níveis (federal, estadual, municipal). Uma nota poderia dirigir-se ao INPI com recomendações sobre como melhorar sistemas de proteção de conhecimentos tradicionais. Outra ao Ministério da Ciência e Tecnologia sobre como integrar salvaguardas para povos tradicionais em políticas de inovação. Outra aos órgãos ambientais sobre como incorporar metodologias de valoração multidimensional em processos de decisão sobre conservação. Este tipo de engajamento com gestores públicos é crucial para que resultados de pesquisa se convertam em mudança sistêmica.

\textbf{Materiais de Comunicação e Engajamento Amplo} amplificarão o alcance para públicos não-acadêmicos. Um \textbf{documentário audiovisual} (20-30 minutos) narrará histórias de comunidades participantes, mostrará na prática como o modelo funciona, incluirá depoimentos de mestres de saberes e gestores sobre utilidade. Este documentário será disponibilizado em plataformas abertas (YouTube, Vimeo) com legendas em português e eventualmente outras línguas. Um \textbf{website dedicado ao projeto} funcionará como portal centralizado de acesso a todos os recursos: documentos, guias, dados, ferramentas interativas. Uma \textbf{aplicação móvel} será desenvolvida tornando componentes do modelo (particularmente a plataforma digital de registro e valoração de conhecimentos) acessível em smartphones, crucialmente importante para comunidades rurais. \textbf{Cartilhas ilustradas} para comunidades explicarão em linguagem muito acessível conceitos como proteção de propriedade intelectual, valoração de conhecimentos, direitos de povos tradicionais.

\subsubsection{Licenciamento Aberto e Replicabilidade}

Um comprometimento central é garantir que o modelo e seus componentes se tornem bem público e passível de replicação por qualquer organização, comunidade ou pesquisador que deseje utilizá-lo. Portanto, o modelo em sua totalidade, seus instrumentos (templates, protocolos, escalas psicométricas), dados de validação (anonimizados conforme consentimento comunitário), e códigos computacionais (algoritmos de machine learning, otimização) serão disponibilizados sob licença aberta — recomendando-se licença Creative Commons BY-SA (Atribuição-CompartilhaIgual) que exige que derivações mantenham a mesma licença aberta, ou para código computacional, licenças GPL ou Apache.

Acompanhando esta disponibilização, será produzida documentação rigorosa de metadados: cada recurso conterá descrição clara de o que é, como foi derivado, que restrições de uso existem (por exemplo, dados sensíveis que podem ser compartilhados apenas sob aprovação ética), como deve ser citado se utilizado em novas pesquisas ou projetos. Esta documentação minimiza fricção para potenciais utilizadores.

Complementando, será estruturado \textbf{programa comprensivo de capacitação} para multiplicadores interessados em adaptar o modelo. Este programa incluirá: workshops presenciais intensivos (3-5 dias) oferecidos em universidades-parceiras sobre operacionalização do modelo; webinários mensais oferecidos gratuitamente abordando temas específicos (Como implementar o módulo de machine learning? Como conduzir validação comunitária? Como interpretar resultados de otimização multiobjetivo?); cursos online auto-paced (autoinstrucionais, em plataformas como Coursera ou Moodle) permitindo aprendizado no próprio ritmo; e acesso a comunidades de prática online (fóruns, grupos de discussão) onde pessoas utilizando o modelo em diferentes contextos compartilham experiências, resolvem problemas, cocriam melhorias.

\subsubsection{Sustentabilidade Institucional de Longo Prazo}

Para que o modelo não seja apenas um artefato de pesquisa que existe enquanto a bolsa financiar o projeto, mas se converta em infraestrutura permanente de valor para comunidades e para o país, será necessário assegurar sustentabilidade institucional. Paralelamente à execução do projeto e validação, o grupo de pesquisa explorará ativamente viabilidades de integração do modelo em estruturas institucionais permanentes.

Uma possibilidade é que componentes do modelo — particularmente a plataforma digital de registro e proteção de conhecimentos tradicionais — sejam integrados como serviço permanente oferecido pelo Núcleo de Inovação Tecnológica (NIT) da UFS. O NIT, por suas atribuições legais, acompanha processos de proteção de propriedade intelectual derivada de pesquisa. Estender este acompanhamento para incluir proteção de conhecimentos tradicionais e uso de ferramentas do modelo seria expansão natural de sua missão.

Alternativamente, órgãos de gestão ambiental (ICMBio, órgão estadual de meio ambiente) poderiam incorporar o modelo em seus processos de planejamento participativo em unidades de conservação. Se comunidades tradicionais que vivem em ou próximo a unidades de conservação utilizarem o modelo para valorar e proteger seus conhecimentos, isso fortaleceria sua participação nas decisões de gestão.

Uma terceira possibilidade é que plataformas colaborativas de economia solidária — redes de cooperativas, mercados justo-solidários, sistemas de comércio comunitário — incorporassem componentes do modelo (particularmente a plataforma de registro de conhecimentos com rastreabilidade via blockchain) para certificação de origem e autenticidade de produtos comunitários.

Em qualquer destas configurações, será crucial: (i) que a comunidade continue tendo controle sobre seus conhecimentos e dados — estruturas de governança devem garantir poder decisório comunitário; (ii) que financiamento seja provido de forma previsível para manutenção operacional, atualizações técnicas e adaptação contínua; e (iii) que feedback contínuo de utilizadores reais seja incorporado em ciclos regulares de melhoria do modelo. Sem estas salvaguardas, a sustentabilidade será ilusória.

O impacto transformador duradouro desta pesquisa não se medirá apenas por publicações, mas por resposta à pergunta: Cinco anos após término do projeto, o modelo continua sendo utilizado por comunidades? Comunidades conseguem autossustentavelmente aplicar o modelo sem depender permanentemente de pesquisadores externos? Políticas públicas incorporaram recomendações derivadas do modelo? Outras comunidades replicaram e adaptaram o modelo com sucesso? Estas são as medidas verdadeiras de sucesso.

